% =====================================================================
%  acknowledgements.tex —— 致谢与免责声明（原书 p.5，位于 Part I 扉页之后、第 1 章之前）
% =====================================================================
\chapter*{致谢与免责声明}
\addcontentsline{toc}{chapter}{致谢与免责声明}
\markboth{致谢与免责声明}{}

我要感谢物理系给我开设这门课程的机会。也希望参与者尽可能多地指出这些讲义中以及课堂上的错误。

一般说来，讲义包含的内容远多于我将在课堂上或能够讲授的内容。同样，讲义中有许多论述可以在不影响理解后续内容的前提下跳过（许多这类非必要的部分标记为“附注”）。类似地，讲义中许多论述只给出数学和/或物理上的事实而不加解释。由于这些原因，读者在讲义中遇到不理解的内容时不必焦虑。我只是希望这套讲义能成为一个起点，让读者由此开始探索代数拓扑的各个主题。但同样由于这套讲义的这种性质，我不希望其中包含错误或不正确的信息。若你发现讲义中有任何错误，请不吝告知我（例如发邮件至 [邮箱隐去]）。

这是我第一次大量借助 AI 准备讲义。更具体地说，我在 \LaTeX{} 写作中使用基于 Chat-GPT 的 GitHub Co-pilot。这个工具在写作过程中建议句子或段落；建议并不总是完全正确，但往往有用，其给出的 \LaTeX{} 宏大多正确。但我担心，从某种意义上说我是在抄袭用于训练 AI 的内容。至少就这套讲义而言，我不认为这有特别的问题，因为课程内容大多是人所共知的、并非我的原创工作。可能唯一原创的部分是内容的选择与呈现的顺序。\footnote{原文注：但这句话本身也是我打出“可能唯一”时 Co-pilot 自动建议的……}若有人注意到任何看起来像是直接抄自其他来源的部分，请告知我。
